%! Author = chtompki
%! Date = 1/14/26
% Example LaTeX document for GP111 - note % sign indicates a comment
\documentclass[12pt]{amsart}
% Default margins are too wide all the way around. I reset them here
\setlength{\hoffset}{-.75in}
\setlength{\textwidth}{6.5in}
\setlength{\voffset}{-.5in}
\setlength{\textheight}{9.0in}
\usepackage{amsmath}
\usepackage{amssymb}
\usepackage{amsthm}
\usepackage{pgfplots}
\usepackage{enumerate}
\usepackage{tikz}
\usetikzlibrary{shadows,arrows,arrows.meta,shapes,positioning,calc}% Required for the Circle and Triangle arrow tips
\usepackage{setspace}
\usepackage{siunitx}

\theoremstyle{conjecture}
\newtheorem{conjecture}{Conjecture}

\theoremstyle{definition}
\newtheorem{definition}{Definition}

\theoremstyle{question}
\newtheorem{question}{Question}

\theoremstyle{theorem}
\newtheorem{theorem}{Theorem}

\title{Seminar Summary:\\April 21, 2026 - Brent Cody\\Connections Between Graph Theory and Music}
\author{Rob Tompkins}
\date{\today}


\begin{document}
    \maketitle
    \date{\today}
    \begin{onehalfspace}
        We start with an introduction to modular synthesizers.
        But we are quickly confronted with a problem about maximal evenness, presenting a 12-cycle $C_{12}$, we are
        curious about the maximal placing of 5 onsets over the cycle of 12.
        We try to space the onsets alternatively at first, and we get a section of three uncovered vertices in the cycle.
        We do a little fiddling, and we end up with the repeating pattern of the black-keys on a piano over the white
        keys on the piano.
        The set $J_{12,5}^0 = \{0,2,4,7,9\}$ is the maximally even subset of $C_{12}$.

        So naturally we ask, what does it mean to be maximally even?
        We come up with the following definition
        \begin{definition}
            (Clough, Douthett, 1991) Let $d(u,v)$ represent the distance from vertex $u$ to vertex $v$ of shortest length where $u,v\in C_n$,
            with $C_n$ being the cycle on $n$-vertices.
            We call a set of vertices $\{a_0,a_1,\ldots, a_{m-1}\}\subseteq C_n$ \emph{maximally even} if whenever there is
            $1\le k \le \lfloor m/2 \rfloor$, we have
            \[
                d(a_i, a_{i+k}) \in \left\{ \left\lfloor \frac{nk}{m} \right\rfloor, \left\lceil \frac{nk}{m}\right\rceil \right\}
            \]
        \end{definition}
        Now we notice something curious, namely that
        \[
            J_{12,5} = \left\{\left\lfloor\frac{12i}{5}\right\rfloor: 0\le i< 5\right\}.
        \]
        So this theorem arises:
        \begin{theorem}
            (Clough, Douthett, 1991) A set $A \subseteq C_n$ with $|A| = m$ is maximally even if and only if it is
            of the form
            \begin{displaymath}
                A = J_{n,m}^r = \left\{ \left\lfloor\frac{ni+r}{m}\right\rfloor: 0 \le i < m\right\}.
            \end{displaymath}
        \end{theorem}

        And it seems that quite naturally the complement of a maximally even set is also maximally even

        \[
            C_n \setminus J_{m,n}^r = J_{n,m-n}^{n-r-1}
        \]
    \end{onehalfspace}
\end{document}
